% --- Parte 1: Introducción y Motivación ---
\section{Introducción y Motivación}
\SectionFrame{Introducción y Motivación}

\begin{frame}{El problema de la inferencia privada}
    \begin{itemize}
        \setlength{\itemsep}{8pt}
        \item \textbf{Ejecución por terceros.} Los modelos de redes neuronales puede que sean ejecutados por terceros, ya sea por modelos propietarios o por facilidad operacional.
        \item \textbf{Los datos son sensibles.} Registros médicos, información financiera o datos personales quedan expuestos a fugas o a un uso indebido.
        \item \textbf{No siempre es posible.} Por cumplimiento legal o por necesidad de confidencialidad, muchos usuarios simplemente no pueden usar estos servicios.
    \end{itemize}
    \vfill
    \begin{center}
        \begin{tikzpicture}[scale=0.9, every node/.style={font=\scriptsize}]
            \node[draw, rounded corners, fill=UCVGrey!30, minimum width=2.2cm, minimum height=1.0cm] (cli) at (0,0) {Cliente};
            \node[draw, rounded corners, fill=UCVGrey!30, minimum width=2.2cm, minimum height=1.0cm] (srv) at (7,0) {Servidor};
            \node[below=1pt of cli, font=\tiny, text=UCVBlue] {datos privados};
            \node[below=1pt of srv, font=\tiny, text=UCVBlue] {modelo entrenado};

            \draw[dashed, gray] (3.5,-0.95) -- (3.5,1.35);
            \node[gray, font=\tiny, anchor=south] at (3.5,1.4) {barrera de confianza};

            \draw[->, thick, UCVOrange] ([yshift=9pt]cli.east) -- node[above, font=\tiny, text=UCVOrange, pos=0.28] {datos en claro} ([yshift=9pt]srv.west);
            \draw[<-, thick, gray!70] ([yshift=-9pt]cli.east) -- node[below, font=\tiny, text=gray, pos=0.28] {resultado} ([yshift=-9pt]srv.west);
        \end{tikzpicture}
    \end{center}
\end{frame}

\begin{frame}{El cifrado homomórfico como solución}
    El cifrado totalmente homomórfico (FHE) permite \textbf{operar sobre datos cifrados sin descifrarlos}. El servidor evalúa el modelo sin ver nunca la información del cliente, y solo este último puede descifrar el resultado.
    \vfill
    \begin{center}
        \begin{tikzpicture}[scale=0.8, every node/.style={font=\scriptsize},
            actor/.style={draw, rounded corners, fill=UCVGrey!30,
                          minimum width=2.1cm, minimum height=1.0cm, align=center},
            dato/.style={draw, rounded corners, fill=UCVBlue!12,
                         minimum width=1.7cm, minimum height=0.6cm, align=center, font=\tiny},
            cifra/.style={draw, rounded corners, fill=UCVOrange!22,
                          minimum width=2.0cm, minimum height=0.6cm, align=center, font=\tiny},
            rotulo/.style={font=\tiny, inner sep=2pt}]

            % Columna del cliente: datos privados arriba, resultado abajo.
            \node[actor] (cli) at (0,0) {\faUser\\ Cliente};
            \node[dato]  (dat) at (0,1.7) {Datos privados};
            \node[dato]  (res) at (0,-1.7) {Resultado};
            \draw[gray!60] (cli.north) -- (dat.south);
            \draw[gray!60] (res.north) -- (cli.south);

            % Textos cifrados: lo unico que llega a cruzar la barrera.
            \node[cifra] (ct1) at (6.2,1.7) {\faLock~ Datos cifrados};
            \node[cifra] (ct2) at (6.2,-1.7) {\faLock~ Resultado cifrado};

            % Servidor de terceros, del otro lado de la barrera.
            \node[actor] (srv) at (10.2,0) {\faServer\\ Servidor};
            \node[below=2pt of srv, rotulo, text=UCVBlue, align=center]
                {evalúa el modelo\\ sobre datos cifrados};

            \draw[dashed, gray] (8.0,-2.6) -- (8.0,2.5);
            \node[gray, rotulo, anchor=south] at (8.0,2.5) {barrera de confianza};

            % Ida: se codifica y cifra antes de cruzar la barrera.
            \draw[->, thick, UCVOrange] (dat.east)
                -- node[rotulo, above, text=UCVOrange] {codificado y cifrado} (ct1.west);
            \draw[->, thick, UCVOrange] (ct1.east) -| (srv.north);

            % Vuelta: sale por el oeste para no pisar el rotulo del servidor.
            \draw[->, thick, gray!70] (srv.west) -- (8.45,0) |- (ct2.east);
            \draw[->, thick, gray!70] (ct2.west)
                -- node[rotulo, below, text=gray] {descifrado y decodificado} (res.east);
        \end{tikzpicture}
    \end{center}
\end{frame}
